\section{Certificate size and comparison with tensor geometry}
\label{sec:complexity50}
\subsection{Finite output bounds}
Here the input to Theorem~\ref{thm:duality49} is explicitly listed.
At a threshold $\delta$, let $v=|V_\delta|$ be the number of essential
factor coordinates, let $n$ be the number of table entries, and let
$m\le2v+2n$ be the number of inequalities in one chamber. Sparsity
below is per chamber. It is not a bound for the union of chambers.
The empty-active case is solved by the zero product and needs no rays.

\begin{proposition}[Explicit certificate and algebraic-output bounds]
\label{prop:complexity50}
For $v\ge1$, write $r_2$ for the rank over $\mathbb F_2$ of the active
sign equations. When these equations are consistent, at most
$2^{v-r_2}$ sign assignments need be tested, and induced-pattern
identifications can only reduce this number. Each chamber has at most
\begin{equation}\label{eq:rays50}
 R(m,v)=\sum_{s=1}^{\min(m,v+1)}\binom{m}{s}
\end{equation}
extreme rays. Each ray has a primitive integer representative with
support at most $v+1$, entries at most $v!$, and sum at most
$h_v=(v+1)!$. A chamber-complete infeasibility certificate thus uses at
most $2^{v-r_2}$ ray/endpoint records, in addition to its input and sign
coverage description. Its ray indices and multiplicities require at
most
\[
 O\bigl(2^{v-r_2}(v+1)(\log(m+1)+v\log(v+1))\bigr)
\]
bits. This bound is on the record data, not a polynomial-time search.

For rational targets with numerator and denominator bit lengths at
most $\tau$, every boundary polynomial on a fixed threshold stratum
has degree at most $h_v$ and can be represented by integer coefficients
of bit length $O(h_v(\tau+1))$. Including the original breakpoints,
there are finitely many such polynomials across at most $n+1$ strata,
with the chamber and ray counts above. Taking the largest $v$ over
the strata gives a degree bound $h_v$ for the optimum threshold.
A minimizing machine has a shared algebraic-radical description:
one isolated real threshold, followed by at most $v$ products of
positive rational powers of its algebraic affine bases. Their common
power denominator is at most $v!$. If one instead insists on a single
number field for all factors, degree at most
$h_v(v!)^v$ suffices for rational input.
\end{proposition}
\begin{proof}
Gaussian elimination gives the sign count. The cone
$\{z\ge0:C^Tz=0\}$ is pointed. On an extreme ray with support $s$, the
restricted columns have rank $s-1$, for otherwise there is a
nontrivial perturbation supported there and the ray is not extreme.
Thus $s\le v+1$. A nonzero cofactor vector of an $(s-1)$-row
submatrix spans the kernel. Entries of $C$ are $0,1,-1$, so the
Leibniz determinant bound gives every cofactor magnitude at most
$(s-1)!\le v!$. Choose its nonnegative sign and divide by the common
integer divisor. Each support supports at most one extreme ray.
This proves \eqref{eq:rays50}, the sum bound, and the record size.
If an infeasible chamber exists, Farkas supplies a nonzero dual
vector; only then is the normalized dual polytope nonempty. Some
extreme ray is strictly negative on the logarithmic right side.
Consequently one violating record per rejected chamber is sufficient,
although finding all the records may enumerate every ray.

On a fixed stratum every base is one, a positive affine rational
function of $\delta$, or its reciprocal. Clearing the positive
denominators in a ray product gives a difference of two products
of affine factors, with total multiplicity at most $h_v$. Its degree
is at most $h_v$. Clear the original rational denominators in the
individual factors before multiplying. The coefficient $\ell_1$ norm
of a product is bounded by the product of the factor norms, so its
coefficient bit lengths are $O(h_v(\tau+1))$. Zero polynomials add no
boundary. Root isolation and separate breakpoint tests give the
stated degree bound, using the largest essential-coordinate count.

At a feasible vertex, a nonsingular $v$-by-$v$ integer row matrix
$A$ gives $y=A^{-1}\log b$. Its determinant has absolute value
$h\le v!$, and its adjugate is integral. Hence every positive factor
$e^{-y_i}$ has its $h$th power in the field generated by $\delta$.
Adjoining these $v$ positive roots enlarges that field by degree at
most $h^v$, which proves the coarse common-field bound. Row
reconstruction is affine in the factors and does not enlarge it.
\end{proof}

A radical witness can be checked without first constructing that large
common field. Check positivity of every base at the isolated threshold,
raise each factor inequality to the positive common power $h$, and
compare the resulting products in the threshold field. The same applies
to each surviving entry constraint. Exact field arithmetic, integer
powers and real-root isolation are sufficient. Their cost can be
exponential in $v$ and grows with the fully expanded polynomial
coefficients; the compressed exponent notation is not a polynomial
bit-complexity claim. With algebraic input, first specify a common
primitive element, its defining polynomial and isolating interval,
and coordinates of each coefficient in that field. The same proof
works over this explicitly represented field; norms or resultants
reduce the threshold equations to rational polynomials. Field-degree
and coefficient-height growth must be charged as well.

For clarity, a feasible set within a fixed sign chamber and threshold
stratum is the intersection of all extreme-ray product conditions
and the surviving upper-endpoint conditions $U_e>0$. This is the
precise interpretation of the final sentence of
Proposition~\ref{prop:tight49}. The full optimizer must also test the
stratum endpoints; no strict-interior assumption is being made.

\subsection{The Segre exponent matrix and the added interval problem}
The incidence vectors in Section~\ref{sec:duality49} are columns of the
Segre exponent matrix, up to transpose convention. An integer relation
between them gives a binomial identity between rank-one tensor entries.
Minimal-support relations are circuits; parity relations describe the
real sign constraints. Zero factors delete complete coordinate slices.
These structures and the distinction between real and complex
completability are treated directly by Kahle--Kubjas--Kummer--Rosen
\cite[Section 2, especially Proposition 2.2, Remarks 2.3--2.5 and
Proposition 2.7]{KKKR17}. In particular, cubic, nonsquarefree tensor
circuits and sign obstructions are already part of that theory.

The present interval problem does not claim those antecedents as new.
Its data are complete bounded response tables with absolute-error
intervals whose endpoints move with $\delta$. A zero is permitted at
an inactive entry, while coordinates forced by active entries must
survive even when they also occur in inactive entries. On that support,
the upper and lower interval inequalities, factor caps and all real
sign chambers are coupled. Farkas' alternative then certifies the
absence of \emph{any} bounded product in that interval box, and the
odd-coordinate compiler turns every feasible product into legal
stochastic rows. A tensor-completion equality on its own neither
provides those moving intervals nor accounts for the charged selector
in a mixture of machines.

\begin{center}\small
\begin{tabular}{p{.20\textwidth}p{.30\textwidth}p{.38\textwidth}}
\toprule
Present result & Closest antecedent & Identical mechanism and precise addition\\
\midrule
Two-state alternative & Gillis--Shitov \cite{GS19}; tensor completion
\cite{KKKR17}; geometric programming \cite{BKVH07}
& Support, signs, exponent relations and logarithms are classical.
The added result couples the complete error box to an exact stochastic
normal form, including zeros and reconstruction.\\[3pt]
Rank-tight simplex theorem & Positive realization cones \cite{BF04}
& Cone invariance is classical. Rank equality removes all hidden suffix
directions at selected cuts, including separated cuts in variable profiles.\\[3pt]
Occupation and four-state frontier & Simplex difference-body identity;
inherited packet budget
& The volume identity is classical. Its repeated application yields an
all-machine narrow-cut budget even when every packet charge is zero.\\[3pt]
Three-epoch cubic optimum & Segre circuits \cite{KKKR17}; the retained cubic example
& The relation is classical. Rational orthogonal, noncommuting commands
realize it with a sharp stochastic optimum exceeding every separate
flattening bound.\\[3pt]
Finite output bounds & Cofactor bounds, Farkas' lemma and real-algebraic isolation
& These tools are classical. The counts price this exact chamber-complete
algorithm and its minimizing-row representation.\\
\bottomrule
\end{tabular}
\end{center}

The general algebraic two-state optimizer is proved but not implemented
in full. The inherited rational threshold checker and rational
bisection are implemented, as is the specialized cubic certificate.
The new executable checks verify the rational permutation construction,
all its words, the flattening witnesses, the simplex volume identities,
and the stated exact numerical frontier bound. They do not implement
an arbitrary higher-width optimizer or constitute independent priority
certification.
